\chapter{智能体大语言模型的交互方式}\label{chap:agents}



\begin{learningbox}
本章将介绍不同的智能体交互方式，以及智能体的记忆。
具体而言，学完本章后，你将能够：
\begin{itemize}[leftmargin=*, nosep]
    \item 识别人工智能发展史上三种不同类型的{\em 智能体（agent）}，以及它们与当今智能体大语言模型的联系。

    \item 识别智能体大语言模型的五种不同{\em 交互方式（interaction style）}。

    \item 理解{\em 记忆（memory）}在智能体大语言模型中的重要性，并识别不同的智能体记忆架构。

\end{itemize}
\end{learningbox}






\section{人工智能中的智能体}
智能体在人工智能中发挥着核心作用。我们首先讨论人工智能领域中智能体的一些背景。

\subsection{智能体在环境中行动与适应}
在人工智能中，能够在世界中行动的实体称为{\em 智能体}。它们对应于在世界中行动的人类。人工智能研究者试图通过智能体这一概念，把握他们所理解的人类行为的本质：通过构建智能来理解智能。

智能体是在环境中行动的自主软件系统。它们 (1) 拥有目标，(2) 感知环境，(3) 作出决策，(4) 采取行动以实现目标；最重要的是，(5) 从经验中学习。
智能体应当根据新信息进行适应，以更好地实现目标；用强化学习的术语来说，就是改进其动作策略（action policy）。

传统计算机程序是确定性的，会以可预测的方式遵循预先设定的指令（这使计算机在从会计到控制等极其重视可预测性的工作中非常有用）。大语言模型智能体的确定性和可预测性则较低。
它们会学习，并不时改变自身行为。它们能够自主推理、规划和执行任务，利用工具为用户完成简单的工作流程，应用领域包括编程、客户服务和数据分析。（它们也会产生幻觉、犯错，而且总体上难以预测、难以依赖。在这一点上，大语言模型智能体更像人类，而不像计算机。）

人工智能领域对智能体的研究由来已久。
区分以下三类智能体很有帮助：(1) 机器人助手智能体；(2) 信念—愿望—意图智能体；(3) 群体智能体。

\subsubsection{机器人助手}
我们首先讨论作为机器人助手的智能体。
这种对人工智能智能体的理解来自工程领域：智能体为人类执行类似人类所做的任务。智能体是一种工具，能够让日常任务更容易完成。例如，工厂里的焊接机器人、飞机上的自动驾驶仪，或者必须在不可预测的环境中寻找安全路线的自动驾驶汽车。

这类智能体的其他例子包括 Alexa 和 Siri 等虚拟助手，帮助软件开发者提高生产力的 CoPilot 等编程助手，以及帮助客户在公司繁杂规则中找到办法、实现投诉或退款等目标的客服机器人。机器人助手智能体可以具有相当强的推理能力 \citep{dorri2018multi}，而大语言模型正越来越多地用于这类助手 \citep{plaat2025agentic}。

\subsubsection{具有信念、愿望和意图的智能体}
另一种对人工智能智能体的理解植根于认知科学。在这里，目标是让智能体对人类认知过程建模，智能体具有信念、愿望和意图（belief–desire–intention，BDI）。认知科学研究人类的思维过程。人工智能智能体可以用来实现认知模型，并开展思维模型实验：模型能否进行逻辑推理，能否在桥牌游戏中与人类玩家合作，能否解构莎士比亚的一部戏剧、创作音乐，或者具有自我意识？认知可以有许多精微的表现形式，例如解释感官输入、引导身体动作，以及与他人共情。

BDI 智能体最初源于逻辑推理传统。它们被赋予目标，并通过行动追求这些目标。这一领域的文献将行动视为实现智能体目标的手段 \citep{rao1995bdi,georgeff1998belief}。

\subsubsection{群体智能体}
第三种对人工智能智能体的理解，强调大量小型智能体之间的交互模式。前面介绍的多智能体系统由少量能够行动和思考的复杂智能体组成：它们可以提供协助，或者具有信念、愿望和意图（BDI）。这一研究路径则关注由大量简单智能体组成的群体。这里强调的是涌现的协作行为，即在超越单个智能体自身推理的规模上涌现的行为。群体能否提供保护，以抵御捕食者？协作发生的前提是什么？一个社会何时和平，何时处于战争状态？什么是超级有机体（super-organism）\citep{johnson2010deconstructing,wilson1989reviving}？

群体智能体（swarm agent）源于具身人工智能的联结主义传统 \citep{brooks1990elephants}。智能体彼此通信，研究者通常关注合作、通信等大规模行为模式如何涌现。蚁群优化（ant colony optimization）是一种著名的群体算法 \citep{dorigo2007ant}，而 Netlogo 是研究这类模式的知名软件环境 \citep{wilensky2015introduction}。\footnote{\url{https://www.netlogo.org}}

多智能体系统和智能体大语言模型适合研究上述三类路径：作为工具的智能体、作为认知实体的智能体，以及帮助我们研究如何实现和平社会（或如何征服另一个社会）的智能体。本书第四部专门讨论作为工具的智能体，第五部则专门讨论认知与智能体社会。

\subsection{作为强化学习智能体的智能体大语言模型}
研究智能体的行为在人工智能领域有着悠久传统。
智能体行为的一个核心方面，是智能体如何在环境中学习。强化学习（reinforcement learning，RL）\citep{sutton2018reinforcement} 和多智能体强化学习 \citep{wong2021multiagent} 研究智能体从反馈中学习的机制。

下一节将讨论智能体大语言模型中的五种交互方式。作为这一概览的准备，我们先用强化学习的语言，为在世界中交互的智能体提供形式化基础。




图~\ref{fig:rl}~展示了一个由智能体和环境组成的强化学习系统。在这一系统中，智能体通过动作与环境交互，环境则返回系统的新状态，以及一个表示新状态优劣的奖励值。

智能体可以从这些奖励值中学习，记住每个状态下哪些动作有利、哪些动作不利。智能体的目标是找到适合该环境的最优行为，也就是每个状态下的最佳动作。智能体利用奖励强化自己对各状态下最佳动作的认识，从而学到一个表示最优行为的函数。这个将各状态映射到动作的函数称为行为策略（behavior policy），记为 $\pi$。最优策略记为 $\pi^*$。
如果智能体通过深度学习进行学习，那么策略就由神经网络的权重 $\theta$ 参数化。

\begin{figure}
  \begin{center}
    \input figures/rl
    \caption{强化学习循环：智能体根据反馈奖励学习参数化策略 $\pi_\theta$；反馈奖励 $r$ 来自环境的转移函数 $T$ 和奖励函数 $R$，对应于智能体的动作 $a$ \citep{sutton2018reinforcement}}\label{fig:rl}
  \end{center}
\end{figure}

环境表示智能体所处的世界。智能体采取动作，改变环境的状态。环境具有一个概率性的转移函数，每当智能体执行一个动作时，该函数就计算新状态。环境还为新状态计算相应的奖励。因此，环境由两个函数建模：转移函数 $T$ 和奖励函数 $R$。

在这一模型中，智能体在环境中行动、接收反馈，并从中学习哪些动作带来了良好的状态转移，哪些没有。正如下节将要看到的，这一模型与智能体大语言模型的概念非常契合。


\begin{figure}
  \begin{center}
    \input figures/rl-llm
    \caption{智能体利用大语言模型将策略 $\pi_{LLM_{P\pi}}$ 参数化的强化学习：智能体依据来自环境的反馈 $r$ 选择动作 $a$}\label{fig:rl-llm}
  \end{center}
\end{figure}

从强化学习的角度看，我们可以把大语言模型用作智能体中的一个函数：$LLM(prompt) \rightarrow answer$。大语言模型可以用于强化学习的不同位置：既可以在智能体中，也可以在环境中。例如，智能体可以用大语言模型来改进策略。它可以用特定提示调用大语言模型，实现策略函数 $\pi_{LLM_{P\pi}}(s,r)$；该函数从环境反馈中学习，在环境可能处于的不同状态下应采取哪些动作（见图~\ref{fig:rl-llm}）。大语言模型得到提示 $P\pi$，被要求利用奖励 $r$ 改进策略，并返回动作 $a$，用于状态 $s$。正如第~\ref{sec:icrl}~节将要介绍的，上下文强化学习（in-context reinforcement learning，ICRL）使用大语言模型来改进策略，而思维树可以使用大语言模型生成后续动作。

\begin{figure}
  \begin{center}
    \input figures/rl-llm2
    \caption{智能体和环境都使用大语言模型实现 $\pi$、$T$ 和 $R$ 的强化学习。每个大语言模型分别得到一个提示 $Px$，以返回动作、状态或奖励值}\label{fig:rl-llm2}
  \end{center}
\end{figure}

大语言模型也可以用于环境中，计算状态的奖励 $LLM_{R(s)}$，或者计算下一个状态 $LLM_{T(s,a)}$，见图~\ref{fig:rl-llm2}。图中，同一个大语言模型通过三个不同提示承担三种不同角色：执行策略函数、转移函数和奖励函数。例如，思维树可以使用大语言模型评估一个状态的奖励。设计智能体大语言模型时有许多选择，下面将继续讨论。我们用强化学习范式来讨论和说明不同的智能体交互方式。


\section{智能体大语言模型中的学习机制}\label{sec:styles}
智能体大语言模型是一个新兴领域。前面各章展示了在智能体场景中使用大语言模型的许多不同方法。本节将回顾本书迄今介绍的方法，并根据提示设计方法与训练方法对它们加以区分。


我们将介绍适用于智能体大语言模型的五种交互方式。
表~\ref{tab:fourstyles}~概括了这五种智能体交互方式。
我们重点关注控制框架在推理阶段如何向大语言模型提供提示，以及模型如何训练。在特定领域中，所有提示方式都可以产生有用的智能体助手和交互伙伴，使它们以符合实际的方式行动。更复杂的提示方式，则是在更具挑战性的领域追求更高性能的过程中发展起来的。

表的第一列给出智能体交互方式的名称，从查询到强化微调。第二列列出各方式的提示类型，从不带任何示例的简单静态指令提示，到动态强化学习提示。第三列展示大语言模型的训练方法：预训练、微调，或两者兼有。第四列给出该交互方式的具体示例（并附文献引用）。

\begin{table}
\footnotesize\centering
\begin{tabular}{l|c|c|c}
{\em 智能体类型} & {\em 提示} & {\em 大语言模型} & {\em 示例} \\ \hline
{\bf A}：查询 & 用户编写的零样本 & 预训练 & GPT-2 \citep{vaswani2017attention}\\
 & 查询或指令 & & \\ \hline
{\bf B}：上下文学习 & 少样本 & 预训练 & GPT-3 \citep{brown2020language}\\
（监督示例） & 示例 & & \\ \hline
{\bf C}：上下文学习 & 思维链步骤 & 预训练 & 思维链 \citep{wei2022chain}\\
（推理步骤） & 示例 & & 零样本 CoT \citep{kojima2022large}\\ \hline
 & 反馈控制框架 & & 思维树 \citep{yao2024tree}\\
{\bf D}：自我反思 & 提示优化 & 预训练 & LLaMea \citep{van2024llamea}\\
 & 搜索或强化学习 & & ReAct \citep{yao2022react}\\ \hline
{\bf E}：强化微调 & 强化学习提示 & 预训练/ & RLVR \citep{lambert2024tulu}\\
 & 优化控制框架 & 微调 & \\
\end{tabular}
    \caption{五种智能体交互方式}\label{tab:fourstyles}
\end{table}





\subsection{智能体交互方式}
下面讨论这些交互方式。在前三种智能体交互方式中，大语言模型的提示是静态的，由用户编写，保持不变。


\subsubsection{A. 查询提示}
第一种交互方式，是通过给出简单指令来与大语言模型交互。图~\ref{fig:rlA}~展示了这种方式。它只包含一条由用户编写的指令。例如，{\em Lily 没赶上火车。她感到很……}，我们希望模型给出{\em 生气}或{\em 难过}之类的回答。这个提示纯粹是一个查询，不包含任何训练示例。其学习方式是零样本学习（zero-shot learning），只有指令。
\begin{figure}
  \begin{center}
    \input figures/rlA
    \caption{方式 A：静态查询}\label{fig:rlA}
  \end{center}
\end{figure}
大语言模型之所以能够正确回答许多这样的查询，是因为它已经在大规模文本语料库上进行过预训练，其中可能包含用户许多潜在问题的答案。模型还经过了针对指令回答的微调。早期大语言模型是这种交互方式的例子，例如 GPT-2 \citep{vaswani2017attention}。


\subsubsection{B. 上下文学习}
接下来讨论上下文学习。
对于机器学习而言，除了监督学习和强化学习，上下文学习是一种令人振奋的新学习方式。大语言模型中的上下文学习通常只需要少量示例，就能取得良好结果，从而回应了机器学习领域的一项主要挑战。

与交互方式 A 相比，上下文少样本学习是一种更复杂的交互方式。
\begin{figure}
  \begin{center}
    \input figures/rlB
    \caption{方式 B：基于少量示例构造提示的上下文监督学习}\label{fig:rlB}
  \end{center}
\end{figure}
当用户在查询之外提供少量示例时，就会发生上下文学习（in-context learning，ICL）。用户以此对解码器模型施加条件，使其激活参数所表示的大量知识中的正确部分。

图~\ref{fig:rlB}~给出了符合这一交互方式的提示示意。文献中采用这种交互方式的一个例子是 GPT-3。介绍上下文学习的论文是 \citet{brown2020language}。一个少样本提示的例子是：{\em 填写省略号处的内容：the sea-la mer，oysters-huitre，the window-la fenetre，cheese-...}，我们希望模型回答 {\em fromage}。

需要注意，只有模型足够大时，上下文学习才会涌现。最初，上下文学习出现在参数量约为 1000 亿的模型中。在这一规模下，模型所含知识足够丰富，提示中的少量示例就能激活模型的正确部分。随着大语言模型方法不断改进，这一规模门槛也在持续降低。

\subsubsection{C. 思维链}
接下来讨论另一种著名方法：思维链学习 \citep{wei2022chain}，见图~\ref{fig:rlC}~上部。正如第~\ref{chap:reasoning}~章所述，开发这一方法，是为了解决大语言模型不擅长求解简单数学文字推理题的问题。原始思维链方法的一个缺点是，每遇到一种新的推理任务，用户都必须编写新的提示。后续发展出了零样本思维链提示，见图~\ref{fig:rlC}~下部。\citet{kojima2022large} 发现，只需在提示中加入一句通用的话：{\em 让我们一步一步地思考}，就已经能够提高许多基准测试上的性能，尽管提升幅度通常不及为特定任务手写的提示。
\begin{figure}
  \begin{center}
    \input figures/rlC
    \caption{方式 C。上：思维链，基于推理步骤示例的大语言模型上下文学习；下：零样本思维链，使用一条通常有效的指令}\label{fig:rlC}
  \end{center}
\end{figure}
与前一种交互方式一样，思维链也是一种上下文少样本学习方法。不过，这时提示中的示例，是模型在给出答案之前应当采取的{\em 步骤}的少量示例，而前一种方式中的示例是事实。这些步骤引导模型沿着一条多步推理链得出正确答案。示例不再是用于激活模型某些部分的答案示例，而是如何得到答案的步骤示例。少样本和零样本思维链的提示示例，已在图~\ref{fig:cot}~中展示，见第~\pageref{fig:cot}~页。

同样，与上下文学习一样，思维链也只在规模足够大的模型中有效；随着我们对大语言模型方法的理解加深，以及模型接受面向思维链的微调，所需规模也在不断减小。此外，通过蒸馏，还可以把大模型学到的推理方法迁移到参数量小几个数量级的模型中 \citep{li-etal-2023-symbolic}。




\subsubsection{D. 自我反思控制框架}
前面介绍的智能体交互方式都只使用一个静态提示，在许多基准测试上取得了令人鼓舞的性能。要在规划问题或其他必须跟踪状态的问题上表现良好，就需要一种能够优化提示、在可能性树中搜索各种选项的推理方式。

下面讨论两种交互方式：它们通过位于大语言模型之外的小型算法，动态调整模型的提示。这个算法称为{\em 控制框架（harness）}。第一种方式是在上下文中进行的（不改变大语言模型参数），第二种方式（下一小节介绍）还会对模型进行微调。

这些交互方式是动态的。
正如上一章关于元学习的讨论所述，动态方法面临的一个挑战，是实现稳定的学习过程。当训练目标不断变化时，学习算法可能发散，需要格外谨慎才能实现收敛。
如何通过自我反思获得可预测的高性能，仍然是一个活跃的研究领域。

我们首先讨论的动态提示设计方法，源于对更复杂推理结构进行建模的需求；这些结构超越了思维链中直线式的步骤序列。最初引入的是简单搜索算法（思维树）。随后，为了管理提示的演化，又引入了更完备的动态提示设计方法（ReAct、reflexion、self refine）。这些方法的控制框架通常引入短期记忆和长期记忆，以便在多次大语言模型调用之间记住状态。

这些方法有趣的一点是，大语言模型承担了不同角色，分别通过不同类型的提示来调用。在思维链任务中，大语言模型最初的角色是{\em 评估}每个推理步骤中的答案。现在，在这种评估提示之外，又加入了另一类提示，用来{\em 生成}新的推理步骤。有时还会加入第三类提示，任务是{\em 改进}提示本身。因此，同一个模型通过不同提示进行自我反思，而这些提示由外部控制框架管理。

\begin{figure}
  \begin{center}
    \input figures/rlD
    \caption{方式 D：利用大语言模型进行自我反思。通过大语言模型的回答在上下文中优化提示：先生成子问题，再评估子问题；这一过程由搜索（上）或强化学习（下）控制}\label{fig:rlD}
  \end{center}
\end{figure}


图~\ref{fig:rlD}~展示了这些思想。上部使用标准树搜索算法控制自我反思，如思维树方法 \citep{yao2024tree}。下部则使用强化学习循环控制自我反思，并利用（状态、动作、奖励）三元组的历史记录更新策略 $\pi$ 的提示（另见第~\ref{sec:icrl}~节）。





\subsubsection{E. 强化微调控制框架}
到目前为止，所有学习都发生在上下文中，模型根据上下文提供的示例给出答案，大语言模型的参数没有改变。尽管 ICL 易于使用，它也限制了学习。如果希望模型更深入地适应某个环境，希望模型记住它在推理阶段学到的内容，那么，允许参数变化就会有所帮助。我们将采用这种方法的研究路径称为{\em 强化微调（reinforcement finetuning）}。RLVR 是一个著名的例子。

图~\ref{fig:rlE}~给出了示意，其中用 $\psi$ 表示可训练的模型参数。注意，它与图~\ref{fig:rlD}~下部的唯一区别，是引入了 $\psi$ 参数。在 ReFT 中，大语言模型参数 $\psi$ 会变化，即根据数据集的奖励进行微调。
应当这样理解这些图：原先推理阶段的学习发生在上下文中；现在，由于允许对参数 $\psi$ 进行微调，这种情况发生了变化。
在传统大语言模型训练流程中，微调是推理之前的一个阶段。然而，在数学推理或编程等许多重要应用中，思维链在推理阶段展开的多步轨迹能够产生准确的奖励，可用于微调。因此，在强化微调中，模型在推理阶段生成训练目标，再用这些目标微调同一个模型的参数，也就是在推理阶段进行微调。强化微调模糊了推理与微调之间的界限，而监督微调则明确区分训练与推理。



\begin{figure}
  \begin{center}
    \input figures/rlE
    \caption{方式 E：对大语言模型参数 $LLM^\psi$ 进行强化微调，同时对策略 $\pi$ 的参数 $\pi_\theta$ 进行强化学习 \citep{luong2024reft}。}\label{fig:rlE}
  \end{center}
\end{figure}

大语言模型训练主要涉及三个训练/推理阶段：(1) 预训练；(2) 微调；(3) 测试（推理）。自我反思（方式 D）在推理阶段进行强化学习。
在方式 E，即强化微调中，我们在微调阶段 (2) 进行强化学习，而不是在测试阶段 (3) 进行。
对于智能体大语言模型而言，这种交互方式使大语言模型能够从智能体与世界的交互中学习。

对于动态智能体交互方式 D 和 E，必须谨慎处理，才能实现稳定收敛。截至本书撰写时，只有少量已发表研究报告了这方面的结果。

\section{智能体大语言模型的记忆架构}\label{sec:memory}
人类（以及其他生命形式）能够记忆。记忆使我们拥有身份认同。我们能记住自己的名字、朋友和家人。记忆使我们学会什么对自己有益、什么有害，也使我们能够调整行为。
当我们因受伤、疾病或衰老而丧失记忆时，人格也会发生变化。严重的记忆丧失会导致身份认同丧失，使我们在世界中漫无目的地游荡。

智能体也需要记忆。本节讨论智能体大语言模型为什么需要记忆，并简要概述为其开发的不同记忆类型。

\subsection{为什么需要记忆}
统计模型 $f(x) \rightarrow \hat{y}$ 会估计输出 $\hat{y}$。这一估计基于关系 $f$，即模型在训练阶段针对输入 $x$ 学到的关系。这种函数式行为没有记忆。任何新信息（训练之后出现的信息）都必须通过输入参数 $x$ 提供给模型。大语言模型也不例外；即使最早的聊天机器人，也会得到特殊的隐藏{\em 系统提示（system prompt）}，将对话历史输入模型，使其能够保持连贯的对话。

随着智能体大语言模型的应用场景不断增多，需要比系统提示方法更先进的技术。
为了更好地理解记忆在智能体大语言模型中的用途，下面讨论记忆的不同应用场景。

\subsubsection{连贯对话}
许多针对大语言模型的机器学习基准测试采用挑战/响应形式，非常适合统计模型的函数式输入/输出行为。然而，大语言模型在现实世界中的主要应用是聊天机器人，人们期望它能与用户进行有意义的对话。为了维持连贯的对话，让对话双方能够对彼此作出反应和回应，聊天机器人必须知道之前说过什么。

因此，智能体大语言模型的首要需求之一，是能够知道对话中之前说过什么。



\subsubsection{适应}
根据第~\ref{chap:intro}~章的定义，我们期望智能体大语言模型能够学习。当机器人在厨房执行切洋葱等任务时，如果刀的位置发生了变化，就应相应调整动作：现在，它应当先从新位置拿起刀，再开始切。智能体与环境交互；如果环境改变，智能体就应当适应。

要么，智能体的预训练策略必须具有很强的通用性，能够适应特定领域中的所有新情境；要么，智能体架构必须足够灵活，能够识别预训练策略何时需要适应、微调，甚至重新训练。为此，它们需要记忆。

正如第~\ref{sec:self-reflection}~节所述，自我反思智能体架构使用的外部提示改进控制框架，需要缓冲区及其他类型的记忆，以显式改进提示。

\subsubsection{身份认同}
经历塑造我们的人格和身份认同。如果大语言模型智能体要建立自己的身份认同，就必须收集经历并保存在记忆中，用于自我反思；或者对经历进行处理，并以其他方式用来调整行为策略。

随着智能体控制框架的记忆架构不断改进，智能体将能够更好地记住自己的经历。届时，我们或许会发现，智能体能够通过自身作为智能体的经历，形成独特的人格。

\subsubsection{持续学习}
传统机器学习针对特定任务，在特定数据集上训练模型。人类的学习则不同，我们的学习在一生中持续演变。
在机器学习中，持续学习（continual learning）研究如何针对一系列新任务训练模型，同时保留先前学到的任务能力 \citep{wang2024comprehensive,li2025improving}。与元学习一样，它研究相关任务之间的泛化。持续学习面临的主要挑战，是克服灾难性遗忘（catastrophic forgetting），以及应对不断变化的数据分布。

持续学习领域提供了支持这类训练的多种架构。



\subsection{记忆架构}
智能体大语言模型领域尚处于早期。直到最近，大语言模型智能体的记忆架构才开始受到关注。下面讨论迄今出现的不同方法。

\subsubsection{微调}
应对记忆需求的一种实用方法，是在新信息出现时就将其纳入模型。我们可以利用模型参数本身就是一种记忆这一事实，在新示例到来时微调模型，而不必开发外部记忆管理算法。

要在智能体场景中实现这一方法，微调必须足够快，在用户察觉延迟之前就完成。LoRA 等参数高效微调方法是一种选择 \citep{yang2024low}。尽管这一方法简洁，不需要外部算法或记忆结构，但对大多数应用而言，微调速度仍然不够快，因此应当开发基于上下文或控制框架的解决方案。

\subsubsection{长上下文窗口与压缩}
第二种方法是管理提示。聊天机器人的标准做法，是在用户文本之外，把经过截断的聊天历史注入提示，提供给大语言模型。这种方法在实践中效果相当好。它的局限在于，上下文窗口的大小有限；而且，注意力机制的时间复杂度为平方级，处理长上下文成本很高。

研究者探索了两类解决方案。首先，开发专门的长上下文架构，以不同方式跟踪注意力 \citep{smith2022simplified}。这是一个活跃的研究领域，\citet{liu2025comprehensive} 对近期方法进行了综述。
其次，开发压缩方案，将提示在语义层面压缩成较短的序列。相关综述见 \citet{li2025prompt}。
提示可以被划分为多个块 \citep{sheng2025dynamic}，以便处理大型提示。一种相关方法使用大语言模型生成代码，递归处理长提示 \citep{zhang2025recursive}。

还可以通过过滤不必要的词元来压缩提示（例如，本句可以压缩为{\em 提示压缩过滤不必要词元}）。
另一种方法是软提示（soft prompt）\citep{lester2021power}。
软提示不再使用离散词元，而是通过反向传播学习嵌入；这些嵌入无需映射到实际的语言词元，就能直接输入 Transformer。这种上下文可以替代纯文本上下文。已有研究报告了可观的压缩比，从十倍 \citep{mu2023learning} 到百倍 \citep{li2025500xcompressor} 不等。




\subsubsection{显式记忆管理}
如果需要在多次大语言模型调用之间存储状态，一种直接的解决方案，是让智能体控制框架在管理提示的同时提供这种存储。第~\ref{sec:self-reflection}~节介绍过，用于自我反思的外部提示改进算法如何加入不同类型的存储，例如 reflexion 方法和思维缓冲方法（图~\ref{fig:reflexion2}~和图~\ref{fig:bot}）。

智能体对存储可能有不同的需求（综述见 \citet{zhang2025survey,chhikara2025mem0}）。我们将讨论这一活跃研究领域中的一些近期研究。提示可以用文本或参数形式存储，相关操作包括读取和写入记忆。此外，压缩和遗忘信息等记忆管理操作也很重要。记忆系统还可以提供额外服务。例如，MemoryBank \citep{zhong2024memorybank} 加入了分层摘要，并使用大语言模型提示，根据用户的交互事件持续构建其人格画像。

可以为智能体加入存储和检索条目的操作。MemGPT \citep{packer2023memgpt} 研究大语言模型如何通过函数调用读写外部记忆（将记忆作为工具）。
A-mem \citep{xu2025mem} 基于邮箱原则，引入了记忆的上下文描述。

显式记忆管理方法的一个缺点是，大语言模型必须显式执行记忆管理，却没有受过如何高效完成这一任务的训练。如果能够预见所有可能的需求，另一种解决方案是手工编写记忆管理代码。
模型上下文协议（Model Context Protocol，MCP）是一组为上下文管理设计的策略。第~\ref{sec:mcp}~节将更深入地介绍这一协议。


\paragraph{学习使用记忆}
Mem1 \citep{zhou2025mem1} 采用另一种方法：每次大语言模型交互之后，智能体都会更新一个紧凑状态，供未来的大语言模型调用共享。所需记忆大小保持恒定，因此状态不会超出上下文窗口。通过在提示中插入推理形式的标签来创建共享状态，并使用强化学习进行训练，以激励状态整合。Mem-$\alpha$ 是另一种使用强化学习训练智能体有效使用记忆的方法 \citep{wang2025mem}。

AgeMem \citep{yu2026agentic} 将记忆管理操作整合到智能体策略中，进行端到端优化。它区分长期记忆与短期记忆，并使用三阶段强化学习方案，训练智能体使用这两类记忆。

这同样是一个活跃的研究领域，预计还会出现新的研究和框架。相关概览见 \citet{zhou2026we}。


\paragraph{向量数据库}
除了把提示存储为人类可读的文本，我们也可以直接存储其词元化后的嵌入。
当需要精确的文本匹配时，以文本形式存储提示有助于检索数据。嵌入编码语义信息，表示哪些词元在含义上彼此接近。因此，如果希望检索语义相近的条目，存储嵌入就很有用。

嵌入可以存储在向量数据库（vector database）中。传统关系数据库以行和列存储数据。向量数据库存储高维数值列表，并提供低延迟访问，因此非常适合大语言模型 \citep{xie2023brief,ma2023comprehensive}。



\section{结语}
智能体是人工智能的一个核心概念。大语言模型显著增强了智能体的通信能力与推理能力。本章讨论了不同类型的智能体，并将它们与人工智能的不同传统联系起来。

智能体会行动。基于前面各章，我们讨论了 5 种主要交互方式。智能体需要记忆，才能与环境进行有意义的交互；本章也讨论了这一新的研究领域。

学习了这么多理论，现在该亲自动手实践了。下一章将介绍一个实践案例研究，让我们在运用前面的一些理论概念方面积累更多经验。




\bibliographystyle{plainnat}
\bibliography{agenticbook,refs,case,synthetic}
